\chapter*{Conclusion Générale}
\addcontentsline{toc}{chapter}{Conclusion Générale}
\vspace{-1.5cm}

Ce projet de fin d’études, réalisé au sein de l’entreprise \textbf{BI NewVision} en partenariat avec l’\textbf{\ac{ONDA}}, s’inscrit dans une dynamique de modernisation des aéroports marocains. Il porte sur le développement d’un \textbf{jumeau numérique} pour optimiser l’affectation des stands avion, en intégrant les contraintes métier et les aléas liés au trafic aérien.


Pour mener à bien ce projet, nous avons commencé par une analyse approfondie du dataset fourni, afin de mieux comprendre les données opérationnelles disponibles et d’en extraire les caractéristiques essentielles pour la modélisation. Cette phase a permis d’identifier des schémas comportementaux et de préparer les données à des fins de simulation et d’optimisation.

Le \ac{SAP} a été modélisé en intégrant plusieurs contraintes réelles telles que la compatibilité avion-stand, les conflits d’occupation et les restrictions d’adjacence. Nous avons défini une fonction objectif multicritère, puis testé différentes approches algorithmiques, exactes et heuristiques, évaluées selon des métriques comme la qualité des affectations, la robustesse et le temps de calcul.

Par la suite, nous avons mis en place une \textbf{architecture microservices}, intégrant un moteur d’optimisation via une API Python, couplé à une interface 3D permettant de simuler les affectations.


Ce projet m’a permis de développer des compétences techniques et méthodologiques variées, notamment en modélisation mathématique, algorithmique, ingénierie logicielle, traitement de données, ainsi qu’en gestion agile de projet au sein d’une équipe interdisciplinaire.

\vspace{1em}

Malgré l’avancement significatif du projet, plusieurs axes d’amélioration restent envisageables pour enrichir et renforcer la solution proposée. Nous envisageons les perspectives suivantes :

\begin{itemize}
	\item Finalisation de l'intégration complète de la solution dans l’environnement client.
	\item Ajout de nouveaux modules fonctionnels comme la gestion des bagages.
	\item Finalisation de l’intégration Cloudera pour l’exploitation directe des données
	\item Enrichissement du dataset pour affiner les performances des algorithmes.
	\item Intégration de contraintes supplémentaires spécifiques aux besoins évolutifs de l’ONDA.
\end{itemize}

Ces perspectives offrent un potentiel réel de développement pour rendre la solution encore plus complète, flexible et adaptable aux futurs besoins du client.

%En conclusion, ce projet de PFE a été une expérience enrichissante et stimulante, nous permettant de développer une application mobile novatrice utilisant la reconnaissance vocale pour guider les touristes à Marrakech. Nous avons relevé plusieurs défis tout au long du processus, tels que le scraping de données pour obtenir des informations précieuses, ainsi que l'adaptation de notre modèle de reconnaissance vocale à plusieurs langues avec de bons résultats.

%L'utilisation de technologies telles que Python, Java et la reconnaissance vocale nous a permis de repousser nos limites et de mieux comprendre les aspects techniques de leur intégration. De plus, nous avons exploré des fonctionnalités hors ligne pour assurer une expérience utilisateur fluide, même sans connexion Internet.  

%\noindent

%Ce projet nous a permis d'acquérir de nouvelles compétences en matière de développement d'applications mobiles avec le machine learning, de travailler avec des technologies émergentes et de résoudre des problèmes complexes. Nous avons également réalisé l'importance de l'encadrement, et nous tenons à exprimer notre profonde gratitude envers nos encadrants pour leur soutien et leurs encouragements tout au long du projet.

%En termes d'impacts potentiels, notre application peut contribuer à améliorer l'expérience touristique à Marrakech en offrant un guide vocal personnalisé et en facilitant la découverte de la ville pour les visiteurs. 

%Grosso modo, ce projet a été une étape importante dans notre parcours universitaire, nous permettant de mettre en pratique nos connaissances et de repousser les frontières de nos compétences. Nous sommes fiers des résultats obtenus.

%Cependant, même avec les fonctionnalités actuelles de notre application, il existe encore plusieurs perspectives intéressantes pour son amélioration. Nous envisageons les perspectives suivantes :
%\begin{itemize}
    %\item Généralisation de l'application pour toutes les villes du Maroc.
    %\item Intégration d'un système de réservation directement dans l'application pour faciliter et centraliser les réservations d'hébergements, de visites guidées, de restaurants, etc.
    %\item Ajout d'une fonctionnalité de chat intégrée permettant aux utilisateurs de contacter directement les responsables des hôtels, des restaurants, et autres établissements pour obtenir des informations, des recommandations ou une assistance personnalisée.
    %\item Intégration d'une carte géographique hors ligne pour permettre aux utilisateurs de naviguer facilement dans la ville de Marrakech sans dépendre d'une connexion Internet constante.
    %\item Mise en place d'un historique des activités pour permettre aux utilisateurs de consulter leurs activités passées, les lieux visités, les commentaires laissés, etc.

%\end{itemize}
 %   Ces perspectives ouvrent la voie à des améliorations futures de notre application et à des fonctionnalités encore plus avancées. Nous sommes impatients de poursuivre notre travail et d'explorer ces possibilités pour offrir aux touristes une expérience de voyage encore plus immersive et pratique à Marrakech.%





















